\documentclass[a4,11pt]{aleph-notas}
% Se puede ver la documentación aquí:
% https://github.com/alephsub0/LaTeX_aleph-notas

% -- Paquetes adicionales
\usepackage{enumitem}
\usepackage{url}
\usepackage{array}
\usepackage{booktabs}
\usepackage{longtable}
\hypersetup{
    urlcolor=blue,
    linkcolor=blue,
}

% -- Datos
\institucion{Facultad de Ciencias Exactas, Naturales y Ambientales}
\carrera{Catálogo STEM}
\asignatura{Álgebra lineal}
\tema{Tema 10: Aplicaciones lineales}
\autor{Andrés Merino}
\fecha{Periodo 2026-01}

\logouno[0.16\textwidth]{Logos/logoPUCE_04_ac}
\definecolor{colortext}{HTML}{1957d1}
\definecolor{colordef}{HTML}{1957d1}
\fuente{montserrat}

\begin{document}

\encabezado 

%%%%%%%%%%%%%%%%%%%%%%%%%%%%%%%%%%%%%%%%
\section*{Resultado de Aprendizaje}
%%%%%%%%%%%%%%%%%%%%%%%%%%%%%%%%%%%%%%%%

%%%%%%%%%%%%%%%%%%%%%%%%%%%%%%%%%%%%%%%%
\subsection*{RdA de la asignatura:}
\begin{itemize}[leftmargin=*]
    \item \textbf{RdA 1:} Comprender los conceptos fundamentales del Álgebra Lineal, incluyendo el estudio de matrices, determinantes, sistemas de ecuaciones lineales y espacios vectoriales, destacando su importancia en el análisis y resolución de problemas matemáticos.
    \item \textbf{RdA 2:} Resolver problemas prácticos y teóricos relacionados con el Álgebra Lineal, utilizando técnicas de cálculo con matrices, determinantes y sistemas de ecuaciones, así como en el análisis de espacios vectoriales.    
    % \item \textbf{RdA 3:} Aplicar los principios y métodos del Álgebra Lineal en una variedad de contextos demostrando comprensión de su relevancia y utilidad en situaciones prácticas y teóricas.
\end{itemize}


%%%%%%%%%%%%%%%%%%%%%%%%%%%%%%%%%%%%%%%%
\subsection*{RdA de la actividad:}
\begin{itemize}[leftmargin=*]
    \item Reconocer y verificar cuándo una función entre espacios vectoriales es una aplicación lineal.
    \item Determinar núcleo, imagen, nulidad y rango de una aplicación lineal.
    \item Analizar inyectividad, sobreyectividad, biyectividad e isomorfismos a partir del núcleo, la imagen y la dimensión.
    \item Construir y utilizar la matriz asociada de una aplicación lineal respecto de bases dadas.
\end{itemize}

%%%%%%%%%%%%%%%%%%%%%%%%%%%%%%%%%%%%%%%%
\section*{Introducción}
%%%%%%%%%%%%%%%%%%%%%%%%%%%%%%%%%%%%%%%%

%%%%%%%%%%%%%%%%%%%%%%%%%%%%%%%%%%%%%%%%
\paragraph{Control de lectura:}
Antes de iniciar la clase, se aplicará el control de lectura del siguiente enlace:
\href{https://gemini.google.com/share/7d22976ac15d}{Control de lectura - Clase 09}.

\paragraph{Pregunta inicial:}
¿Cómo se reconoce una aplicación lineal y qué información esencial aporta su núcleo o su matriz asociada?

%%%%%%%%%%%%%%%%%%%%%%%%%%%%%%%%%%%%%%%%
\section*{Desarrollo}
%%%%%%%%%%%%%%%%%%%%%%%%%%%%%%%%%%%%%%%%

%%%%%%%%%%%%%%%%%%%%%%%%%%%%%%%%%%%%%%%%
\subsection*{Actividad 1: Aplicaciones lineales, núcleo e imagen}

La clase inicia con una recuperación breve de la definición de aplicación lineal y de sus propiedades básicas. Luego, mediante clase magistral y práctica guiada, se estudian el núcleo, la imagen, la nulidad, el rango y los criterios de inyectividad y sobreyectividad. Se cierra con la construcción de la matriz asociada a una aplicación lineal y su uso para representar transformaciones entre bases.

\paragraph{¿Cómo lo haremos?}
\begin{itemize}[leftmargin=*]
    \item \textbf{Clase magistral:} se presentan las definiciones de aplicación lineal, núcleo, imagen, nulidad, rango, isomorfismo y matriz asociada. Se usan ejemplos en $\mathbb{R}^n$, matrices y polinomios. Se usa el resumen \href{https://fcena-puce.github.io/AlgebraLineal-05-N0048/01Resumen/Resumen10.pdf}{Resumen10.pdf}.
    \item \textbf{Resolución de ejercicios:} se guiará a los estudiantes en la solución de ejercicios del documento \href{https://fcena-puce.github.io/AlgebraLineal-05-N0048/04Ejercicios/Ejercicios10.pdf}{Ejercicios10.pdf}, identificando aplicaciones lineales, calculando núcleos e imágenes, verificando isomorfismos y construyendo matrices asociadas.
    \item \textbf{Visualización de video:} se recomendará el estudio con el video: \href{https://youtu.be/kYB8IZa5AuE?si=51z3YxrVrpOQwNTV}{Transformaciones lineales y matrices.}.
\end{itemize}

\paragraph{Verificación de aprendizaje:}
\begin{enumerate}[leftmargin=*]
    \item ¿Qué condiciones debe verificar una función para ser lineal?
    \item ¿Cómo se determina el núcleo y la imagen de una aplicación lineal?
    \item ¿Cómo se decide si una aplicación lineal es inyectiva, sobreyectiva o un isomorfismo?
    \item ¿Cómo se construye la matriz asociada de una aplicación lineal respecto de bases dadas?
\end{enumerate}

%%%%%%%%%%%%%%%%%%%%%%%%%%%%%%%%%%%%%%%%
\section*{Cierre}
%%%%%%%%%%%%%%%%%%%%%%%%%%%%%%%%%%%%%%%%

\paragraph{Tarea:}
Resolver del libro \href{https://catalogobiblioteca.puce.edu.ec/cgi-bin/koha/opac-detail.pl?biblionumber=86081}{Fundamentos de álgebra lineal de Larson}: sección 6.1, ejercicios: 1, 3, 5, 7, 9, 11, 17; sección 6.2, ejercicios: 1, 3, 5, 7; sección 6.3, ejercicios: 1, 3, 5, 7, 9, 23, 25.

\paragraph{Pregunta de investigación:}
\begin{enumerate}[leftmargin=*]
    \item ¿Por qué el núcleo determina la inyectividad de una aplicación lineal?
    \item ¿Qué información geométrica codifica la imagen de una transformación lineal?
    \item ¿Cómo cambian las matrices asociadas al cambiar de base?
\end{enumerate}

\paragraph{Para el próximo tema:}  
Visualizar el video \href{https://youtu.be/YJfS4_m_0Z8?si=kFS74G3kXhWGWq4L}{Transformaciones lineales}.
Realizar la actividad de clase invertida: \href{https://fcena-puce.github.io/AlgebraLineal-05-N0048/07Act-ClaseInvertida/03ClsInv-AplicacionLineal.pdf}{03ClsInv-AplicacionLineal.pdf}.



\end{document}
